
\documentclass{article}
\usepackage{colortbl}
\usepackage{makecell}
\usepackage{multirow}
\usepackage{supertabular}

\begin{document}

\newcounter{utterance}

\centering \large Interaction Transcript for game `codenames', experiment `concrete', episode 9 with reflexion{-}Qwen3.5{-}27B.
\vspace{24pt}

{ \footnotesize  \setcounter{utterance}{1}
\setlength{\tabcolsep}{0pt}
\begin{supertabular}{c@{$\;$}|p{.15\linewidth}@{}p{.15\linewidth}p{.15\linewidth}p{.15\linewidth}p{.15\linewidth}p{.15\linewidth}}
    \# & $\;$A & \multicolumn{4}{c}{Game Master} & $\;\:$B\\
    \hline

    \theutterance \stepcounter{utterance}  
    & & & \multicolumn{2}{p{0.3\linewidth}}{
        \cellcolor[rgb]{0.9,0.9,0.9}{
            \makecell[{{p{\linewidth}}}]{
                \texttt{\tiny{[GM$|$GM]}}
                {'hidden': {'team': ['panther'], 'innocent': ['cranberry', 'dustpan', 'longbow', 'toothbrush', 'burger'], 'opponent': [], 'assassin': []}, 'revealed': {'team': {'team': [], 'innocent': [], 'opponent': [], 'assassin': []}, 'opponent': {'team': [], 'innocent': [], 'opponent': [], 'assassin': []}}}
            }
        }
    }
    & & \\ \\

    \theutterance \stepcounter{utterance}  
    & & \multicolumn{4}{p{0.6\linewidth}}{
        \cellcolor[rgb]{0.9,0.9,0.9}{
            \makecell[{{p{\linewidth}}}]{
                \texttt{\tiny{[P1$\langle$GM]}}
                \texttt{You have attempted this same task before and did not succeed. These are the plans you wrote after those attempts.} \\
\texttt{Trial \#1: New plan: The previous attempt was aborted because the clue "sting" is semantically related to the kill word "eagle" (as in a sting operation or potentially confusing associations, but more critically, the game rules often imply that if a clue is too broad or ambiguous, it risks triggering a loss, though the specific error here was likely that "sting" is not a valid clue for "toad" or "barracuda" in a way that satisfies the strict semantic requirement without ambiguity, or perhaps the system flagged "sting" as containing a part of a word or being too close to a forbidden concept. However, looking closer at the rules: "The clue word... cannot be one of the words in the lists or contain parts of them." "Sting" does not contain parts of the target words. The most likely reason for the ABORT is that "sting" is a valid word for "wasp" but a very weak or invalid association for "barracuda" and "toad" in the context of the game's strict validation, or the system considers "sting" to be a "kill" word association (e.g., "sting" {-}> "eagle" via "stingray" confusion? No, that's a stretch).} \\
\\ 
\texttt{Actually, re{-}reading the error: "ABORT —— the Game Master rejected one of your messages because it violated the game rules". The most common rule violation in this specific game format (Codenames) when an ABORT happens is often a formatting issue or a clue that is a substring of a target word. "Sting" is not a substring. Wait, "sting" is not in the lists.} \\
\\ 
\texttt{Let's reconsider the "kill" word: "eagle". Is "sting" related to "eagle"? Not directly.} \\
\texttt{Is it possible the error was that I included "toad" and "barracuda" which are not strongly related to "sting"? The rules say "The clue word has to be semantically related to the target words". If the validator is strict, "sting" is not a valid clue for "toad" (toads don't sting) or "barracuda" (they bite, they don't sting). This might be the violation: claiming a relation that doesn't exist.} \\
\\ 
\texttt{However, the prompt says "Your plan must focus on how to produce a well{-}formed and legal response."} \\
\texttt{Let's look at the example 2 again. The error was a 3{-}letter word.} \\
\texttt{In my case, the format `CLUE: sting` and `TARGETS: wasp, barracuda, toad` looks correct.} \\
\texttt{Perhaps the issue is that "sting" is a substring of "stingray" which isn't there, or maybe the system considers "sting" to be a "kill" word? No} \\
\\ 
\texttt{Trial \#2: New plan: My previous strategy failed because I prioritized grouping multiple words under broad clues (like "animal" for "toad" and "panther") which inadvertently triggered guesses of distractor words (like "chimpanzee" or "warthog") or opponent words, ending my turn prematurely. Additionally, I left "egg" unguessed until the end, likely because I didn't find a safe, specific clue for it that didn't overlap with other categories. Next time, I will prioritize high{-}confidence, specific clues that link to only one or two words to minimize the risk of triggering distractors. Specifically, I will use "shell" or "yolk" for "egg" early on to secure it without ambiguity. For "toad" and "panther," instead of the broad "animal," I will use "frog" (specific to toad) or "big cat" (specific to panther) to avoid the "chimpanzee" trap. I will also avoid clues like "bite" or "jump" that have strong associations with distractors (warthog, chimpanzee) and focus on unique attributes of the target words to ensure my partner guesses only my team's words.} \\
\\ 
\texttt{Trial \#3: New plan: My previous strategy failed because I provided clues for only one word at a time ("frog" for "toad", "yolk" for "egg", "cat" for "panther"), causing my turn to end prematurely after each single guess. In this game, the goal is to guess *all* team words in a single turn to win quickly; ending the turn after one guess allowed the opponent to play and eventually win. Next time, I will analyze the remaining team words to find a single clue that semantically connects multiple targets simultaneously. For the set "egg, toad, panther", I should have used a broader but safe clue like "animal" (connecting toad and panther) or "hatch" (connecting egg and toad) to list multiple targets in the `TARGETS` field, thereby maximizing the number of words guessed before the turn ends. I will prioritize grouping words over safety of single{-}word clues to prevent the opponent from getting a turn.} \\
\\ 
\\ 
\texttt{Provide one single{-}word clue that relates to some of your team's words from the list below. You can choose to how many and to which words your clue relates to.} \\
\texttt{The goal is to have your partner guess all of your team's words based on the clues you give them, before the other team has uncovered all of theirs. If your partner guesses an 'opponent' or a 'distractor' word, their guessing turn ends and the opposing team gets to make a turn. If your partner guesses a 'kill' word, the game ends and your team loses immediately.} \\
\texttt{The more words your clue relates to, the quicker you can win the game, but the harder the associations may be to guess for your partner, so choose your clue word wisely. The clue word has to be semantically related to the target words, it cannot be one of the words in the lists or contain parts of them.} \\
\\ 
\texttt{Always give your single{-}word clue and your comma{-}separated list of related target words in the following format and make your answers as short as possible, never include any other text than is required in this form:} \\
\\ 
\texttt{CLUE: <WORD>} \\
\texttt{TARGETS: <TARGETS>} \\
\\ 
\texttt{Your team words are: panther.} \\
\texttt{Your opponent's team words are: .} \\
\texttt{Distractor words are: cranberry, dustpan, longbow, toothbrush, burger.} \\
\texttt{The kill words are: .} \\
            }
        }
    }
    & \\ \\

    \theutterance \stepcounter{utterance}  
    & \multicolumn{4}{p{0.6\linewidth}}{
        \cellcolor[rgb]{0.9,0.9,0.9}{
            \makecell[{{p{\linewidth}}}]{
                \texttt{\tiny{[P1$\rangle$GM]}}
                \texttt{CLUE: cat} \\
\texttt{TARGETS: panther} \\
            }
        }
    }
    & & \\ \\

    \theutterance \stepcounter{utterance}  
    & & & \multicolumn{2}{p{0.3\linewidth}}{
        \cellcolor[rgb]{0.9,0.9,0.9}{
            \makecell[{{p{\linewidth}}}]{
                \texttt{\tiny{[GM$|$GM]}}
                \texttt{cat} \\
            }
        }
    }
    & & \\ \\

    \theutterance \stepcounter{utterance}  
    & & & \multicolumn{2}{p{0.3\linewidth}}{
        \cellcolor[rgb]{0.9,0.9,0.9}{
            \makecell[{{p{\linewidth}}}]{
                \texttt{\tiny{[GM$|$GM]}}
                ['panther']
            }
        }
    }
    & & \\ \\

    \theutterance \stepcounter{utterance}  
    & & & \multicolumn{2}{p{0.3\linewidth}}{
        \cellcolor[rgb]{0.9,0.9,0.9}{
            \makecell[{{p{\linewidth}}}]{
                \texttt{\tiny{[GM$|$GM]}}
                {'word': 'panther', 'assignment': 'team'}
            }
        }
    }
    & & \\ \\

    \theutterance \stepcounter{utterance}  
    & & & \multicolumn{2}{p{0.3\linewidth}}{
        \cellcolor[rgb]{0.9,0.9,0.9}{
            \makecell[{{p{\linewidth}}}]{
                \texttt{\tiny{[GM$|$GM]}}
                \texttt{opponent has won} \\
            }
        }
    }
    & & \\ \\

    \theutterance \stepcounter{utterance}  
    & \multicolumn{4}{p{0.6\linewidth}}{
        \cellcolor[rgb]{0.95,0.95,0.95}{
            \makecell[{{p{\linewidth}}}]{
                \texttt{\tiny{[P2$\rightarrow$P1]}}
                \texttt{Post{-}episode Reflexion (not sent to GM):} \\
\texttt{New plan: My previous attempt failed because I provided a clue ("cat") that, while semantically correct for "panther," was too broad and likely triggered a guess of a distractor word (e.g., "dustpan" via the "pan" suffix or a general animal association) or simply ended the turn prematurely without securing a win condition in a multi{-}word scenario. Although "panther" was the only target, the strategy of using a generic category word is risky. Next time, I will select a more specific, unique attribute or a less common synonym that strongly points to the target while minimizing overlap with distractors. For "panther," instead of "cat," I will use a specific descriptor like "black" (referencing the black panther) or "leopard" (a specific type of panther) to ensure the association is tight and less likely to be confused with "dustpan" or other distractors. If multiple targets were present, I would prioritize a clue that links them uniquely without touching distractor concepts.} \\
            }
        }
    }
    & & \\ \\

\end{supertabular}
}

\end{document}
